\chapter{D\_python\_workflow.py}
\label{ch:pythonWorkflow}

This is the script for the Python motion compensation workflow. The motion compensation is done by applying a rigid body transformation (rotation + time-interpolated translation) to all points in an input \texttt{.las} file.

\section{How It Works}
This script corrects the motion-induced distortions by:
\begin{enumerate}
    \item Loading the distorted \texttt{.las} file.
    \item Loading the converted Pulsalyzer format trajectory file and removing duplicate timestamps for linear interpolation.
    \item Verifying that the \texttt{.las} file timestamps are within the input trajectory time range, if not, stop execution.
    \item The trajectory is filtered to the corresponding time window of the input \texttt{.las} file.
    \item Checking trajectory linearity with the condition $R^2 \ge 0.999$. If condition fails, stop execution.
    \item Linear interpolation functions are then built for $x$, $y$, and $z$ as a function of time.
    \item Building a rotation matrix from user-defined boresight Euler angles (intrinsic XYZ).
    \item Processing the \texttt{.las} file in chunks to manage memory.
    \item Rotating each point in a chunk using the rotation matrix and adding the linearly interpolated position from the linear interpolation functions.
    \item Each chunk is written incrementally to an output \texttt{.las} file.
\end{enumerate}

\section{User-Configurable Parameters}
\begin{table}[!ht]
    \centering
    \begin{tabularx}{\textwidth}{llX}
        \toprule
        Variable                            & Default & Description                                                                                \\ \midrule
        input\_las\_file                    & R"..."  & Path to distorted .las file (e.g.: \texttt{\detokenize{R"C:\Folder\file.las"}})            \\
        input\_pulsalyzer\_trajectory\_file & R"..."  & Path to converted trajectory .txt file (e.g.: \texttt{\detokenize{R"C:\Folder\file.txt"}}) \\
        x\_rot, y\_rot, z\_rot              & 0       & Boresight rotation angles in degrees                                                       \\
        chunk\_size                         & 1000000 & Points processed per iteration                                                             \\ \bottomrule
    \end{tabularx}
\end{table}

\section{Rotation Convention}
The rotation matrix is built using the \textbf{Intrinsic XYZ} rotation convention.

The module \texttt{scipy.spatial.transform} of the Python library \texttt{scipy} is used for building the rotation matrix. The method used for building the rotation matrix from input angles in degrees following the intrinsic XYZ rotation convention is:

\begin{lstlisting}[numbers=none, language=python]
Rotation.from_euler('XYZ', [x_rot, y_rot, z_rot], degrees=True).as_matrix()
\end{lstlisting}

For more information on the usage of the module and rotation convention notations, refer to the documentation available at:

\url{https://docs.scipy.org/doc/scipy-1.16.2/reference/generated/scipy.spatial.transform.Rotation.html}

The boresight angles (obtained from the Boresight Calibration (\autoref{ch:scriptE})) used in the thesis for the Python workflow were:

\begin{equation*}
    x\_rot = \qty{2.5}{\degree}, \quad y\_rot = \qty{184}{\degree}, \quad z\_rot = \qty{48}{\degree}
\end{equation*}

\section{Output}

The transformed point cloud is saved as a \texttt{.las} file to the same directory as the input \texttt{.las} file using the following naming convention:

\begin{verbatim}
python_transformed_{x_rot}_{y_rot}_{z_rot}_{input_filename}.las
\end{verbatim}

The output file retains all scalar fields of the input file. The LAS header offset is updated to the mean position of the filtered trajectory.